\documentclass[11pt]{article}

\usepackage[margin=1in]{geometry}
\usepackage{amsmath,amssymb,amsthm,mathtools}
\usepackage{hyperref}

\newcommand{\EE}{\mathbb{E}}
\newcommand{\PP}{\mathbb{P}}
\newcommand{\Var}{\mathrm{Var}}
\newcommand{\cF}{\mathcal{F}}
\newcommand{\R}{\mathbb{R}}

\title{IE 622 --- Practice Problem Set 13\\Detailed Solutions}
\author{}
\date{}

\begin{document}

\maketitle

\noindent
This note contains complete solutions to all problems in Practice Problem Set 13, together
with the main concepts used in each solution.

\section*{Problem 1}

\textbf{Statement.} Consider a DTMC $\{X_n\}$ and suppose state $N$ is such that
$p_{N,N}=0$. Let $g(i)$ denote the probability that this chain eventually enters state $N$
given that it starts in state $i$. Show that $\{g(X_n)\}$ is a martingale with respect to
$\{\sigma(X_1,\dots,X_n)\}$.

\subsection*{Concepts used}

\begin{itemize}
\item The Markov property.
\item First-step analysis for hitting probabilities.
\item A function $g$ satisfying $g(i)=\sum_j p_{ij}g(j)$ is called \emph{harmonic} for the chain.
\item If $g$ is harmonic, then $g(X_n)$ is a martingale.
\end{itemize}

\subsection*{Solution}

To make the martingale statement correct at every time, interpret $g(i)$ as the probability
of hitting $N$ at some \emph{strictly future} time when the chain is currently at state $i$.
Equivalently, if the chain is started at state $i$ at the present time, then
\[
g(i)=\PP_i(\tau_N^+<\infty),
\qquad
\tau_N^+ := \inf\{m\ge 1 : X_m=N\}.
\]
With this interpretation, the assumption $p_{N,N}=0$ causes no difficulty even when the
current state is already $N$.

Let
\[
\cF_n := \sigma(X_1,\dots,X_n).
\]
We first show that $g$ satisfies the first-step relation
\[
g(i)=\sum_j p_{ij}g(j)
\qquad \text{for every state } i.
\]
Indeed, condition on the next state:
\[
g(i)
= \sum_j \PP_i(X_{1}=j)\,\PP_i(\tau_N^+<\infty \mid X_1=j)
= \sum_j p_{ij} g(j).
\]
The last equality uses the Markov property and time-homogeneity: once the chain moves to $j$,
the probability of ever hitting $N$ in the future is exactly $g(j)$.

Now compute the conditional expectation of $g(X_{n+1})$:
\begin{align*}
\EE[g(X_{n+1}) \mid \cF_n]
&= \EE\!\left[\EE[g(X_{n+1}) \mid X_n] \,\middle|\, \cF_n\right] \\
&= \EE\!\left[\sum_j p_{X_n j} g(j) \,\middle|\, \cF_n\right] \\
&= g(X_n).
\end{align*}
Thus
\[
\EE[g(X_{n+1}) \mid \cF_n] = g(X_n),
\]
which is exactly the martingale property. Therefore $\{g(X_n)\}$ is a martingale with
respect to $\{\cF_n\}$.

\section*{Problem 2}

\textbf{Statement.} Give an example of i.i.d.\ integrable random variables
$\{X_1,X_2,\dots\}$ and a random time $N$ such that the conclusion of Wald's lemma fails:
\[
\EE\left[\sum_{n=1}^N X_n\right] \ne \EE[N]\EE[X_1].
\]

\subsection*{Concepts used}

\begin{itemize}
\item Wald's lemma requires more than just integrability; a key hypothesis is that $N$ must be
      a suitable stopping time (or one must impose an equivalent condition).
\item If $N$ depends on \emph{future} values of the sequence, Wald's lemma can fail.
\end{itemize}

\subsection*{Solution}

Take $\{X_n\}$ to be i.i.d.\ Rademacher random variables:
\[
\PP(X_n=1)=\PP(X_n=-1)=\frac12.
\]
Then each $X_n$ is integrable and
\[
\EE[X_1]=0.
\]

Now define the random time
\[
N=
\begin{cases}
1, & \text{if } X_2=1,\\
2, & \text{if } X_2=-1.
\end{cases}
\]
Clearly $N$ is integrable, and
\[
\EE[N]=\frac12\cdot 1+\frac12\cdot 2=\frac32.
\]

But
\[
\EE[N]\EE[X_1] = \frac32 \cdot 0 = 0.
\]

Now compute the left-hand side:
\[
\sum_{n=1}^N X_n
=
\begin{cases}
X_1, & \text{if } X_2=1,\\
X_1+X_2, & \text{if } X_2=-1.
\end{cases}
\]
Hence
\begin{align*}
\EE\left[\sum_{n=1}^N X_n\right]
&= \EE[X_1] + \EE[\mathbf{1}_{\{X_2=-1\}}X_2] \\
&= 0 + \left(-1\right)\PP(X_2=-1) \\
&= -\frac12.
\end{align*}
Therefore
\[
\EE\left[\sum_{n=1}^N X_n\right] = -\frac12 \ne 0 = \EE[N]\EE[X_1].
\]

So this gives the required counterexample.

\medskip
\noindent
\textbf{Why Wald's lemma fails here.} The random time $N$ is not a stopping time with respect
to the natural filtration of $\{X_n\}$, because the event $\{N=1\}=\{X_2=1\}$ depends on the
future value $X_2$.

\section*{Problem 3}

\textbf{Statement.} Let $X$ be a random variable with $-a \le X \le b$ and $\EE[X]=0$. If
$f$ is a convex function, show that
\[
\EE[f(X)] \le \frac{1}{a+b}\bigl(bf(-a)+af(b)\bigr).
\]

\subsection*{Concepts used}

\begin{itemize}
\item A convex function lies below every chord joining two points on its graph.
\item Since $X \in [-a,b]$, each value of $X$ can be written as a convex combination of $-a$
      and $b$.
\item The condition $\EE[X]=0$ simplifies the coefficients after taking expectation.
\end{itemize}

\subsection*{Solution}

Fix $x \in [-a,b]$. Then $x$ can be written as
\[
x = \lambda(-a) + (1-\lambda)b,
\qquad
\lambda := \frac{b-x}{a+b}.
\]
Indeed, since $x \in [-a,b]$, we have $0 \le \lambda \le 1$.

Because $f$ is convex,
\[
f(x) \le \lambda f(-a) + (1-\lambda)f(b).
\]
Substituting the value of $\lambda$ gives
\[
f(x) \le \frac{b-x}{a+b}f(-a) + \frac{a+x}{a+b}f(b).
\]
Now apply this inequality to the random variable $X$ and take expectations:
\begin{align*}
\EE[f(X)]
&\le \EE\left[\frac{b-X}{a+b}f(-a) + \frac{a+X}{a+b}f(b)\right] \\
&= \frac{1}{a+b}\left(\bigl(b-\EE[X]\bigr)f(-a) + \bigl(a+\EE[X]\bigr)f(b)\right).
\end{align*}
Since $\EE[X]=0$, this becomes
\[
\EE[f(X)] \le \frac{1}{a+b}\bigl(bf(-a)+af(b)\bigr).
\]
This is exactly the desired inequality.

\section*{Problem 4}

\textbf{Statement.} Let $\{X_n\}$ be the Galton--Watson branching process discussed in class.
If $\mu=1$ and $\PP(Z_1^1=1)<1$, then show that the population dies out almost surely.

\subsection*{Concepts used}

\begin{itemize}
\item A Galton--Watson process is determined by its offspring distribution.
\item If $Z$ denotes the number of children of one individual, then the probability generating
      function is
      \[
      \phi(s)=\EE[s^Z], \qquad 0\le s\le 1.
      \]
\item The extinction probability $q$ is the smallest fixed point of $\phi$ in $[0,1]$.
\item In the critical but non-degenerate case ($\mu=\phi'(1)=1$ and $\PP(Z=1)<1$), the only
      fixed point is $1$.
\end{itemize}

\subsection*{Solution}

Let $Z$ be the offspring random variable and let
\[
\phi(s)=\EE[s^Z], \qquad 0\le s\le 1.
\]
The mean number of offspring is
\[
\mu=\phi'(1)=1.
\]

Let $q$ be the extinction probability when the process starts with one individual. A standard
fact about Galton--Watson processes is that $q$ is the smallest solution in $[0,1]$ of
\[
\phi(s)=s.
\]
So it is enough to show that the only fixed point of $\phi$ in $[0,1]$ is $1$.

Now define
\[
h(s)=\phi(s)-s.
\]
Then
\[
h(1)=\phi(1)-1=0,
\qquad
h'(1)=\phi'(1)-1=\mu-1=0.
\]

We now show that $h(s)>0$ for every $0\le s<1$.

Because $\PP(Z=1)<1$ and $\EE[Z]=1$, the offspring distribution cannot be concentrated at
$1$. Therefore there must be some mass at values $\ge 2$; otherwise, if $Z$ only took the
values $0$ and $1$, its mean would be strictly less than $1$ unless $\PP(Z=1)=1$.
Hence
\[
\PP(Z\ge 2)>0.
\]
Therefore, for $0<s<1$,
\[
\phi''(s)=\EE[Z(Z-1)s^{Z-2}] > 0.
\]
So $\phi$ is strictly convex on $(0,1)$, and therefore $h(s)=\phi(s)-s$ is also strictly
convex on $(0,1)$.

For a strictly convex function, the graph lies strictly above each tangent line except at the
point of tangency. Since the tangent line to $h$ at $s=1$ is
\[
h(1)+h'(1)(s-1)=0,
\]
we get
\[
h(s)>0 \qquad \text{for every } 0\le s<1.
\]
In other words,
\[
\phi(s)>s \qquad \text{for every } 0\le s<1.
\]
Thus there is no fixed point of $\phi$ in $[0,1)$, so the only fixed point in $[0,1]$ is
$s=1$.

Hence the extinction probability is
\[
q=1.
\]
Therefore the Galton--Watson process dies out almost surely.

\section*{Problem 5}

\textbf{Statement.} Let $\{X_1,\dots,X_n\}$ be independent random vectors uniformly distributed
on the circle of radius $1$ centered at the origin. Let $T=T(X_1,\dots,X_n)$ denote the
length of the shortest path connecting these $n$ points. Show that
\[
\PP\bigl(|T-\EE[T]|\ge a\bigr)\le 2\exp\{-a^2/32\}.
\]

\subsection*{Concepts used}

\begin{itemize}
\item The key fact is that $T$ is a bounded random variable.
\item For any random variable $Y$ satisfying $m\le Y\le M$ almost surely, Hoeffding's lemma
      implies
      \[
      \PP(|Y-\EE[Y]|\ge a)\le 2\exp\!\left(-\frac{2a^2}{(M-m)^2}\right).
      \]
\item So the job is simply to find a deterministic upper bound on $T$.
\end{itemize}

\subsection*{Solution}

First note that
\[
T\ge 0.
\]
So we only need an upper bound.

Arrange the points in cyclic order around the circle. If we connect consecutive points in that
order and then delete one edge, we obtain a path visiting all $n$ points. Since $T$ is the
length of the \emph{shortest} path connecting the points, it is no larger than the length of
this particular path.

Now the length of the chord joining two consecutive points is at most the length of the
corresponding arc of the unit circle. Summing over all consecutive pairs around the full
circle, the total arc length is exactly $2\pi$. Hence the polygonal cycle through the points
has total length at most $2\pi$, and after deleting one edge we still have a path of length at
most $2\pi$. Therefore
\[
0\le T\le 2\pi < 8.
\]

Now apply Hoeffding's inequality for a bounded random variable with $m=0$ and $M=2\pi$:
\[
\PP\bigl(|T-\EE[T]|\ge a\bigr)
\le
2\exp\!\left(-\frac{2a^2}{(2\pi)^2}\right).
\]
Since $2\pi<8$, we have
\[
\frac{2}{(2\pi)^2} > \frac{1}{32}.
\]
Therefore
\[
2\exp\!\left(-\frac{2a^2}{(2\pi)^2}\right)
\le
2\exp\!\left(-\frac{a^2}{32}\right).
\]
Thus
\[
\PP\bigl(|T-\EE[T]|\ge a\bigr)\le 2\exp\{-a^2/32\}.
\]

\section*{Problem 6}

\textbf{Statement.} For binary vectors $x$ and $y$ of length $n$, let
\[
d_H(x,y)=\sum_{i=1}^n |x_i-y_i|
\]
denote the Hamming distance between $x$ and $y$. Let $A$ be a finite set of vectors, and let
$X=(X_1,\dots,X_n)$, where the $X_i$ are i.i.d.\ Bernoulli$(1/2)$ random variables. Let
\[
D=\min_{y\in A} d_H(X,y),
\qquad
\mu=\EE[D].
\]
For $a>\mu$, find an upper bound for $\PP(D\ge a)$.

\subsection*{Concepts used}

\begin{itemize}
\item The function
      \[
      f(x_1,\dots,x_n)=\min_{y\in A} d_H(x,y)
      \]
      satisfies a bounded-difference property.
\item Changing one coordinate of $x$ can change $f(x)$ by at most $1$.
\item Therefore McDiarmid's inequality applies.
\end{itemize}

\subsection*{Solution}

Define
\[
f(x_1,\dots,x_n)=\min_{y\in A} d_H(x,y).
\]
Then $D=f(X_1,\dots,X_n)$.

Suppose two binary vectors $x$ and $x'$ differ in exactly one coordinate. For every fixed
$y\in A$, changing one coordinate changes $d_H(x,y)$ by at most $1$. Hence
\[
|d_H(x,y)-d_H(x',y)|\le 1
\qquad \text{for every } y\in A.
\]
Taking the minimum over $y\in A$ preserves this bound, so
\[
|f(x)-f(x')|\le 1.
\]
Thus $f$ satisfies the bounded-difference condition with constants
\[
c_1=\cdots=c_n=1.
\]

By McDiarmid's inequality, for every $t>0$,
\[
\PP(D-\EE[D]\ge t)\le \exp\!\left(-\frac{2t^2}{\sum_{i=1}^n c_i^2}\right)
=\exp\!\left(-\frac{2t^2}{n}\right).
\]
Now set
\[
t=a-\mu >0.
\]
Then
\[
\PP(D\ge a)=\PP(D-\mu\ge a-\mu)\le \exp\!\left(-\frac{2(a-\mu)^2}{n}\right).
\]

So a valid upper bound is
\[
\boxed{\PP(D\ge a)\le \exp\!\left(-\frac{2(a-\mu)^2}{n}\right), \qquad a>\mu.}
\]

\section*{Problem 7}

\textbf{Statement.} Consider a sequence of random variables $\{X_n\}$ defined as follows.
Let $X_1=1$. Given $X_1,\dots,X_{n-1}$, define $X_n$ to be a Poisson random variable with
parameter $X_{n-1}$. What can you say about this sequence when $n$ is large?

\subsection*{Concepts used}

\begin{itemize}
\item Conditional on $X_{n-1}=k$, a Poisson$(k)$ random variable can be written as the sum of
      $k$ independent Poisson$(1)$ random variables.
\item Hence $\{X_n\}$ is a Galton--Watson branching process with offspring distribution
      Poisson$(1)$.
\item This is a critical branching process because the mean offspring number is $1$.
\item By Problem 4, a critical non-degenerate Galton--Watson process dies out almost surely.
\end{itemize}

\subsection*{Solution}

Condition on $X_{n-1}=k$. Then
\[
X_n \mid (X_{n-1}=k) \sim \mathrm{Poisson}(k).
\]
But a Poisson$(k)$ random variable can be represented as
\[
Y_1+\cdots+Y_k,
\]
where $Y_1,\dots,Y_k$ are i.i.d.\ Poisson$(1)$. Therefore each of the $k$ current individuals
produces an independent number of children with distribution Poisson$(1)$.

So $\{X_n\}$ is exactly a Galton--Watson branching process started from one ancestor with
offspring distribution
\[
Z \sim \mathrm{Poisson}(1).
\]
Its mean is
\[
\mu=\EE[Z]=1,
\]
and
\[
\PP(Z=1)=e^{-1}<1,
\]
so the process is critical and non-degenerate.

By Problem 4, the extinction probability is $1$. Hence
\[
X_n \to 0 \qquad \text{almost surely}.
\]
In fact, once the process hits $0$, it stays at $0$ forever, because a Poisson$(0)$ random
variable is identically $0$.

Thus, for large $n$, the population is zero with probability tending to $1$:
\[
\PP(X_n=0)\to 1.
\]

\medskip
\noindent
\textbf{Important remark.} Although $X_n\to 0$ almost surely, the mean stays constant:
\[
\EE[X_n \mid X_{n-1}] = X_{n-1}
\quad \Longrightarrow \quad
\EE[X_n]=\EE[X_{n-1}]=\cdots=\EE[X_1]=1.
\]
So this is a classic critical branching-process phenomenon: almost sure extinction, but with
rare surviving realizations that can still keep the expectation equal to $1$.

\end{document}
